% ID: SEL-D0043-OA-E04
% ESTADO: cotejada. Segunda pasada visual y borrador compilado revisados.
% Fuente: Selectividad/sel_2014_matematicas_aplicadas/EXAMEN -ANDAL-2.pdf
% SHA-256: fc025a25faa338cbb9d8d288443f0322dfbcc0464f7474a16a84e8dfe47f81d1
% Página: 1; opción A; ejercicio original 4.

Con el fin de estudiar el precio medio del litro de gasolina en una provincia en un determinado día, se seleccionan al azar ese día 9 estaciones de servicio y se observan los siguientes precios, en euros, de un litro de gasolina:
\[
1.3,\quad1.2,\quad1.4,\quad1.27,\quad1.25,\quad1.32,\quad1.37,\quad1.38,\quad1.23.
\]
Se sabe que el precio del litro de gasolina se distribuye según una ley Normal con desviación típica igual a 0.18 euros.
\begin{enumerate}[label=\alph*)]
\item \textbf{(1.5 puntos)} Obtenga un intervalo de confianza, al 95\%, para estimar el precio medio del litro de gasolina.
\item \textbf{(1 punto)} Calcule el tamaño muestral mínimo necesario para estimar el precio medio del litro de gasolina con un error no superior a 0.08 euros, con el mismo nivel de confianza.
\end{enumerate}
